\chapter{Технические методы}
\begin{important}
Не стоит забывать, что следующие методы необходимо использовать, по возможности, совместно друг с другом и с методами, перечисленными в предыдущей главе.
\end{important}
\section{Прокси}

\section{VPN}

\section{SSH-туннели}

\section{Настойка браузера Firefox}
\subsection{Подмена информации, передаваемой браузером}
\subsubsection{User-agent}
\subsubsection{Часовой пояс}
\subsubsection{Предпочитаемые языки}
\subsection{Дополнения}
\subsubsection{HTTPS Everywhere}
\subsubsection{BetterPrivacy}
\subsubsection{Ghostery}
\subsubsection{TrackerBlock}
\subsubsection{Beef Taco}
\subsubsection{User Agent Changer}

\section{Анонимные оверлейные сети}
\textbf{Оверлейные сети} --- это такие сети, которые работают поверх другой уже работающей сети.
\subsection{Tor}
\begin{wrapfigure}[8]{r}{0.25\linewidth}
\includegraphics[width=\linewidth]{Tor}
\caption{Логотип Tor}
\end{wrapfigure}
\textbf{Tor} --- анонимная оверлейная сеть, использующая принцип луковой маршрутизации, исходные коды которой распространяются на условиях лицензии BSD\cite{tor_license}.\\
\textbf{Луковая маршрутизация} --- технология анонимного обмена информацией, использующая многократное шифрование и пересылку через цепочки узлов. Каждый луковый маршрутизатор в цепочке удаляет слой шифрования и пересылает сообщение дальше, согласно полученным инструкциям, где все повторится. И так до тех пор, пока сообщение не достигнет адресата. Такое название технология получила из-за сходства данного процесса с очисткой луковицы.\\
\subsubsection{Установка}
\subsubsection{Использование}
\subsubsection{Недостатки}
\begin{important}
Не забывайте, что при использовании обычного Интернета через Tor, последняя нода в цепочке (exit-нода) видит трафик нешифрованным.
\end{important}
\subsection{I2P}
\begin{wrapfigure}[6]{r}{0.25\linewidth}
\includegraphics[width=\linewidth]{I2P}
\caption{Логотип I2P}
\end{wrapfigure}
\textbf{I2P} --- анонимная оверлейная сеть, использующая принцип чесночной маршрутизации, исходные коды которой распространяются на условиях нескольких свободных лицензий\cite{i2p_license}. В отличии от Tor, который в первую очередь направлен на доступ к сайтам обычного интернета (хотя в нем и существуют скрытые сервисы, аналогичные ипсайтам в I2P, а в I2P можно получить доступ к внешнему Интернету, используя аутпрокси), основной целью I2P является доступ именно к скрытым ресурсам --- ипсайтам. Ипсайт от обычного вебсайта отличает только его нахождение в сети I2P.\\
\textbf{Чесночная маршрутизация} --- вариант луковой, отличающийся тем, что несколько <<луковиц>> пересылаются совместно, что усложняет установку авторства сообщений.
\subsubsection{Установка}
\subsubsection{Использование}
\subsubsection{Недостатки}
\subsection{Freenet}
\begin{wrapfigure}[9]{r}{0.25\linewidth}
\includegraphics[width=\linewidth]{Freenet}
\caption{Логотип Freenet}
\end{wrapfigure}
\textbf{Freenet} --- анонимная оверлейная сеть, использующая принципы P2P и F2F и распространяющаяся на условиях GNU GPL v2\cite{freenet_license}. В отличии от Tor и I2P, позволяющих размещать любые ресурсы внутри сети, Freenet по сути представляет собой хранилище статичных данных.\\
\textbf{P2P (Peer-to-peer)} --- компьютерная сеть, в которой все участники равны и выполняют одновременно роль как клиента, так и сервера.\\
\textbf{F2F (Friend-to-Friend)} --- разновидность P2P-сети, в которой все соединения разрешаются только с доверенными узлами.
\subsubsection{Установка}
\subsubsection{Использование}
\subsubsection{Недостатки}

\section{Анонимные платежи}
\subsection{Анонимные пластиковые карточки}
\subsection{Bitcoin}
\subsection{Прочие платежные системы}

\section{IM-сервисы}
\subsection{I2P-Messenger}
\subsection{TorChat}
\subsection{Cryptocat}

\section{Ремейлеры}
\subsection{Типы ремейлеров}
\subsection{Mixmaster}
\subsection{Mixminion}

\section{Прием почты}
\subsection{I2P-Mail}
\subsection{I2P-Bote}
\subsection{TorMail}
\subsection{Privacybox}

\section{Шифрование данных}
\subsection{Truecrypt}
\subsection{dm-crypt}
\subsection{GPG}

\section{Шифрование передаваемой информации}
\subsection{Off-the-Record Messaging (OTR)}
\subsection{GPG и Jabber}
\subsection{GPG м любая другая передаваемая информация}

\section{Стеганография}
\subsection{steghide}
\subsection{OpenStego}

\section{Альтернативные DNS}
\subsection{Namecoin}
\subsection{OpenNIC}
\subsection{Собственный кеширующий DNS сервер}
